% BANKS-SOURCE: pdf=42 print=32 kind=section id=3.7
\subsection{The scattering matrix and the LSZ formula}

% BANKS-SOURCE: pdf=42 print=32 kind=prose
Most of the real experimental data to which QFT can be applied are the results of
scattering experiments. It is an experimental fact that there exist (approximately) stable
single-particle states in the world, as well as states of multiple particles at large relative
space-like distances, which behave, to a very good approximation, like free particles.
Interactions fall off at large spatial separation, and the \emph{cluster property} of QFT provides
a neat mathematical explanation for this \cite{banks-ref-16,banks-ref-17}. The cluster property states that, at
large space-like separation, the connected parts of Green functions fall to zero. In
perturbation theory, this follows from the falloff of a single Feynman propagator. If
all particles are massive, the falloff is exponential, whereas if there is no mass gap we
expect power-law falloff. In this case, a variety of behaviors is possible, and there is not
always a scattering theory.

The idea of scattering theory is to derive formulae for idealized amplitudes in which
some number of widely separated particles come in from past infinity, interact, and
go off to become a state of (possibly different) space-like separated particles in the
asymptotic future. The central quantity one computes is the \emph{scattering} or $S$-matrix: the
amplitude

% BANKS-SOURCE: pdf=42 print=32 kind=display
\[
\langle \mathrm{out}\,f_1\ldots f_n\mid g_1\ldots g_m\,\mathrm{in}\rangle
\]

for $m$ particles with asymptotic wave functions $g_i$ to turn into $n$ particles with asymptotic
wave functions $f_j$. The $g_i$ and $f_j$ run over complete sets of single free-particle states.

Non-relativistic quantum mechanics has a well-developed scattering theory \cite{banks-ref-18,banks-ref-19}.
In that theory one proves that the S-matrix for potential scattering is a \emph{partial isometry},
a unitary transformation on a subspace (the scattering states) of the Hilbert space
orthogonal to all bound states of the particles. The \emph{in} and \emph{out} states are two different
bases for this subspace and the S-matrix is the transformation relating them. In QFT
all bound states can be simply thought of as additional particles. So the entire Hilbert
space will be spanned by scattering states, as long as we treat bound states as separate
particles.

The basic assumption of scattering theory is thus that the Hilbert space has two
complete bases of states, each of which is isomorphic to the Fock space of some col-
lection of free particles. A list of the masses and spins of the particles completely
specifies the spectrum of the Hamiltonian. We have noted in Problem 2.3 how diffi-
cult it is to discuss the eigenstates of an interacting quantum field theory, and claimed
that the introduction of Green functions side-steps all these difficulties. The Lehmann--
Symanzik--Zimmermann (LSZ) formula \cite{banks-ref-20} shows us how to find the masses and spins

% BANKS-SOURCE: pdf=43 print=33 kind=prose
of all single-particle states, and to compute the S-matrix, in terms of Green functions.
To derive it, we assume the existence of a single-particle state of mass $m$ and spin $0$\footnote[10]{The extension to other spins is straightforward and we will not go into the details, but simply use the appropriate formulae when necessary.}.
We also assume the existence of a local field $\Phi(x)$, such that

% BANKS-SOURCE: pdf=43 print=33 kind=display
\[
\langle0|\Phi(x)|p\rangle=\sqrt{\frac{Z}{(2\pi)^{D-1} 2\omega_{\mathbf p}}}\ee^{+\ii p\mathbin{\cdot}x},
\]

with $Z\neq0$. The functional form of the matrix element follows from Lorentz invariance
(prove it if you don't believe me). $\Phi$ might be proportional to the Lagrangian field of
an interacting field theory, or it might be a complicated function of such fields. We will
let $Z[J]$ denote the generating functional of Green functions of $\Phi$.
\BanksImplicitHook{I-CH03-016}
\BanksImplicitHook{I-CH03-015}

Consider a source of the following form: $J=\sum_iJ^i_{\rm in}+\sum_jJ^j_{\rm out}$. Each component of
this composite source is defined as follows:

% BANKS-SOURCE: pdf=43 print=33 kind=display
\[
\int\dd^D x\,J^i_{\rm in}(x)\Phi(x)=\frac{\ii}{\sqrt Z}\epsilon^i_{\rm in}\lim_{t\to\infty}\int\dd^{D-1}\mathbf x\left(\Phi\,\partial_0\phi^{i*}_{\rm in}-\phi^{i*}_{\rm in}\,\partial_0\Phi\right).
\]

There is an analogous formula for the outgoing source. Here $\epsilon^i_{\rm in}$ is infinitesimal and
all formulae are to be interpreted by expansion to lowest order in all the $\epsilon^i_{\rm in/out}$. $\phi^i_{\rm in}$
is a normalizable, positive-energy solution of the KG equation. If $\Phi$ were a free field
of mass $m$, this formula would just define the creation operator for a single-particle
state with wave function $\phi^i_{\rm in}$, the spatial integral would be time-independent and the
limit superfluous. For any local field $\Phi$, the source creates a state localized around
an infinitely distant spatial point as $t\to\infty$. Furthermore, if we consider the matrix
element of $\langle\eta|\int\dd^D x\,J(x)\Phi(x)|\Psi\rangle$ between states $|\Psi\rangle$ and $|\eta\rangle$ of fixed $D$-momentum, then, unless
$(P^\mu_\Psi-P^\mu_\eta)^2=-m^2$, the limit will vanish. If we choose all of the incoming and outgoing
wave functions\footnote[11]{The outgoing wave functions are negative-energy solutions of the KG equation, and we take the limit $t\to-\infty$, instead of $+\infty$.} to be localized around different asymptotic directions as $t\to\pm\infty$,
then this operator acts just like an \emph{in} creation operator even in the interacting theory\footnote[12]{This is the \emph{assumption} that particle interactions fall off at large distances. It can be motivated/proven using certain plausible axioms about QFT \cite{banks-ref-21,banks-ref-22}.}.
Similarly, the \emph{out} part of the source acts like a sum of annihilation operators. Since
$Z[J]=\langle0|T\exp(+\ii\int\dd^D x\,J(x)\Phi(x))|0\rangle$, all the annihilation operators sit to the left of all of the
creation operators. Thus, it is plausible that the generating functional with this source is
nothing but the generating functional for the scattering matrix. That is, the coefficient of
$\epsilon^1_{\rm in}\ldots\epsilon^m_{\rm in}\epsilon^1_{\rm out}\ldots\epsilon^n_{\rm out}$ in the expansion of $Z[J]$ for this source\footnote[13]{We emphasize that other terms in the expansion of the composite source have no particular use. They correspond to creating or annihilating multiple particles in exactly the same scattering state.} is the scattering-matrix
element

% BANKS-SOURCE: pdf=43 print=33 kind=display
\[
\langle \mathrm{out}\,f_1\ldots f_n\mid g_1\ldots g_m\,\mathrm{in}\rangle.
\]

% BANKS-SOURCE: pdf=44 print=34 kind=prose
where the single-particle states $f_i$ and $g_j$ are defined by

% BANKS-SOURCE: pdf=44 print=34 kind=display
\[
\phi^j_{\rm in}=\int\frac{\dd^{D-1}\mathbf p}{\sqrt{2\omega_{\mathbf p}(2\pi)^{D-1}}}\ee^{+\ii p\mathbin{\cdot}x}g_j(\mathbf p)
\]

and

% BANKS-SOURCE: pdf=44 print=34 kind=display
\[
\phi^i_{\rm out}=\int\frac{\dd^{D-1}\mathbf p}{\sqrt{2\omega_{\mathbf p}(2\pi)^{D-1}}}\ee^{-\ii p\mathbin{\cdot}x}f_i(\mathbf p).
\]

Conventional time-dependent perturbation theory offers an alternative definition
of the scattering matrix as the infinite-$T$ limit of the Dirac-picture evolution operator
$U_D(T,-T)$, and it can be verified in all examples that this coincides with our definition
in terms of $Z[J]$. More importantly, the LSZ formula can be massaged into a form
that allows us to find particle masses and scattering amplitudes directly from Green
functions. We first write

% BANKS-SOURCE: pdf=44 print=34 kind=display
\[
\int\dd^D x\,J^i_{\rm in}(x)\Phi(x)=\ii\epsilon^i_{\rm in}\int\dd^D x\,\partial_0\left[\Phi\,\partial_0\phi^i_{\rm in}-\phi^i_{\rm in}\,\partial_0\Phi\right].
\]

This is not quite right. The boundary term at $t=-\infty$ gives what we want, but the
term at $t=\infty$ contains a creation operator for a state localized in the future. However,
if all of the final states are orthogonal to all of the initial states, this term just acts to
the left on the vacuum, and vanishes. In the case where this is not true (the case of
partially forward scattering) the LSZ formula we are about to write must be slightly
modified. In practice, this modification is never of any consequence. Part of the forward-
scattering amplitude is due to disconnected processes happening independently at large
space-like separation. The rest can be obtained by analytic continuation from non-
forward scattering. We now act with the outer time derivative on the KG scalar product,
and exploit its Wronskian form and the fact that $\phi^i_{\rm in}$ satisfies the KG equation to
write

% BANKS-SOURCE: pdf=44 print=34 kind=display
\[
\int\dd^D x\,J^i_{\rm in}(x)\Phi(x)=\ii\int\dd^D x\left[(\nabla^2-m^2)\phi^i_{\rm in}\Phi-\partial_0^2\Phi\,\phi^i_{\rm in}\right]
=\ii\int\dd^D x\,\phi^i_{\rm in}[\Box-m^2]\Phi.
\]

We can do a similar manipulation for all of the in and out sources. For outgoing
particles, we want to pick up the creation operator rather than the annihilation operator,
so we complex conjugate the source equation, and end up with the complex conjugate
of the positive-energy outgoing wave function. The result is the LSZ formula for the
$S$-matrix

% BANKS-SOURCE: pdf=44 print=34 kind=display
\[
\begin{aligned}
\langle \mathrm{out}\,f_1\ldots f_n\mid g_1\ldots g_m\,\mathrm{in}\rangle
={}&\left(\frac{\ii}{\sqrt Z}\right)^{m+n}
\int\prod_k\dd^D x_k\,\phi^k_{\rm in}(x_k)\prod_j\dd^D y_j\,(\phi^j_{\rm out})^*(y_j)\\
&\quad\times\prod_k[\Box_{x_k}-m^2]\prod_j[\Box_{y_j}-m^2]\\
&\quad\times\langle0|T\Phi(x_1)\ldots\Phi(x_m)\Phi(y_1)\ldots\Phi(y_n)|0\rangle.
\end{aligned}
\]

% BANKS-SOURCE: pdf=45 print=35 kind=prose
The product of differential operators includes one KG operator for each integration
variable. It is conventional in scattering theory to replace the normalizable initial- and
final-state wave packets by plane waves. As we shall see, this causes some trivial infinities
to show up in cross sections, but they are easily dealt with. Recall again that the field
$\Phi$ in the LSZ formula \emph{might} be an elementary field, $\phi$ in the Lagrangian, but need not
be. All that is required is that it have a finite amplitude to create single-particle states,
which is verified by finding a pole in its two-point function.

The LSZ formula has two remarkable properties. The first is that it is symmet-
ric between in and out states except for the fact that the out wave functions are
negative-energy solutions and the in wave functions positive-energy solutions of the KG
equation. This leads one to surmise that different scattering amplitudes (e.g. electron--
electron scattering and electron--positron annihilation) are just analytic continuations
of each other. These \emph{crossing symmetry} relations were discussed in a heuristic manner
in the introduction. The hard part in proving them is the proof of analyticity \cite{banks-ref-23}.

Secondly, in momentum space, the formula reads

% BANKS-SOURCE: pdf=45 print=35 kind=display
\[
\begin{aligned}
\langle \mathrm{out}\,q_1\ldots q_n\mid p_1\ldots p_m\,\mathrm{in}\rangle
={}&\left(\frac{\ii}{\sqrt Z}\right)^{m+n}
\prod_i\left[\sqrt{\frac{1}{2\omega_{\mathbf p_i}(2\pi)^{D-1}}(p_i^2+m^2)}\right]\\
&\quad\times\prod_j\left[\sqrt{\frac{1}{2\omega_{\mathbf q_j}(2\pi)^{D-1}}(q_j^2+m^2)}\right]\\
&\quad\times G_{n+m}(-q_1\ldots-q_n;p_1\ldots p_m).
\end{aligned}
\]

$G_{n+m}$ is the Fourier transform of the time-ordered product of $n+m$ $\Phi$ fields. Note
that for the outgoing states $q_i$ is the physical, positive-energy, $D$-momentum, but the
Fourier transform is evaluated at negative-energy outgoing momenta. Because of trans-
lation invariance, the Fourier transform contains a delta function of the sum of all its
arguments, and because of the previous remark, this is just energy--momentum con-
servation $\delta^D(\sum p_i-\sum q_j)$. When we square the amplitude to get a cross section, one
of these momentum-space delta functions gives us $\delta^D(0)$, which is interpreted as the
space-time volume. In a translation-invariant system, we want to compute the prob-
ability per unit volume per unit time for a plane wave to interact. With normalizable
wave packets, this infinity does not occur.

The most remarkable part of this formula is that all the momenta are on mass shell
and the formula contains what appears to be a product of zeros. One concludes that
the S-matrix element will be zero unless the Green function has a single pole at the
mass shell on each external leg. The connection between full connected Green func-
tions and 1PI Green functions shows how easy it is for this to happen. The full Green
functions are evaluated as tree diagrams whose vertices are 1PI functions and whose
internal and external lines are full propagators. In momentum space, each external leg
has a factor $W_2(p)$, the full connected two-point function. The Lehmann representa-
tion shows that each of these Green functions has a simple pole, giving precisely the
multiple-pole structure that we need, in order for the LSZ formula to predict a finite
$S$-matrix.
